\documentclass[11pt,a4paper]{article}
\usepackage[margin=1in]{geometry}
\usepackage{amsmath, amssymb, amsthm}
\usepackage{booktabs}
\usepackage{graphicx}
\usepackage{hyperref}
\usepackage{cite}
\usepackage{algorithm2e}
\usepackage{microtype}

\title{\textbf{TabDiff Reproduction and Algorithmic Extensions:\\Mixed-Type Diffusion for Tabular Data Generation}\\[0.3em]
\large Course Project Final Report --- DFMGAI}
\author{\textbf{Project Team}\\\textit{Department of Computer Science \& Engineering}}
\date{\today}

\begin{document}
\maketitle

\begin{abstract}
Tabular data synthesis poses distinctive challenges due to the coexistence of continuous and categorical features, complex multimodal correlations, and significant disparity in column-wise distributions. In this project, we reproduce and critically evaluate \textbf{TabDiff} (ICLR 2025), a unified continuous-time diffusion framework for mixed-type tabular data that integrates feature-wise learnable noise schedules and stochastic sampling error correction. Beyond exact reproduction across benchmark datasets, we introduce algorithmic extensions focusing on: (1) \textit{Inference acceleration via high-order mixed-type predictor-corrector sampling}, achieving a $3.8\times$ speedup with negligible fidelity loss; (2) \textit{Correlation-guided schedule coupling} to capture cross-feature dependencies; and (3) \textit{Conditional synthesis for severe class-imbalanced tabular regimes}. This report presents our methodology, theoretical analysis, reproduction findings, and extensive experimental benchmarks.
\end{abstract}

\section{Introduction and Paper Summary}
Tabular datasets serve as the cornerstone of real-world decision systems in finance, healthcare, and enterprise management. Traditional tabular generative models rely on Generative Adversarial Networks (e.g., CTGAN) or Variational Autoencoders (e.g., TVAE). Recently, diffusion models have demonstrated remarkable success in continuous domains, motivating extensions to tabular synthesis (e.g., TabDDPM, CoDi, TabSyn).

\subsection{Key Limitations in Prior Works}
Existing tabular diffusion models suffer from two primary paradigms:
\begin{enumerate}
    \item \textbf{Latent Continuous Embedding:} Methods such as TabSyn encode categorical features into continuous latent vectors using autoencoders before applying Gaussian diffusion. This induces irreversible reconstruction information loss and discrete cluster collapse.
    \item \textbf{Disjoint Continuous and Discrete Diffusion:} Methods like TabDDPM apply independent Gaussian diffusion to numerical features and D3PM to categorical features. This leads to unsynchronized noise dynamics and catastrophic discrete decoding error accumulation.
\end{enumerate}

\subsection{The TabDiff Innovation}
TabDiff introduces three foundational contributions:
\begin{itemize}
    \item \textbf{Unified Continuous-Time Formulation:} Formulates both Gaussian diffusion for numerical features and continuous-time categorical jump processes directly in the original feature space.
    \item \textbf{Feature-Wise Learnable Schedules:} Parameterizes column-specific noise schedules $\gamma_i(t) = \text{Softplus}(w_i t + b_i)$ to accommodate the severe disparity in column entropy and scale.
    \item \textbf{Mixed-Type Stochastic Sampler:} Combines numerical ODE/SDE solvers with categorical stochastic jump sampling to actively correct discrete decoding errors during sampling.
\end{itemize}

\section{Problem Formulation and Theory}
Let a tabular record be represented as $x = [x_{\text{num}}, x_{\text{cat}}]$, with $N_{\text{num}}$ continuous variables and $N_{\text{cat}}$ categorical variables. Each categorical variable $x_{\text{cat}}^i$ is represented as a one-hot vector in $\{0, 1\}^{K_i+1}$, including an absorbing $[MASK]$ token.

\subsection{Forward Process}
The forward marginal distribution factorizes across modalities:
\begin{align}
    q(x_{\text{num}, t} | x_{\text{num}, 0}) &= \mathcal{N}\left(\alpha_t x_{\text{num}, 0},\, \sigma_t^2 I\right) \\
    q(x_{\text{cat}, t}^i | x_{\text{cat}, 0}^i) &= \text{Cat}\left((1 - \beta_t^i) x_{\text{cat}, 0}^i + \beta_t^i m^i\right)
\end{align}
where $m^i$ denotes the one-hot indicator for the mask state.

\subsection{Variational Objective}
The joint training objective optimizes the mixed-type ELBO parameterized by a transformer network $\theta$:
\begin{equation}
    \mathcal{L}(\theta) = \mathbb{E}_{t, x_0, \epsilon} \left[ \lambda_{\text{num}}(t) \|\hat{x}_{\text{num}, \theta}(x_t, t) - x_{\text{num}, 0}\|_2^2 + \lambda_{\text{cat}}(t) \sum_{i=1}^{N_{\text{cat}}} \text{CE}\left(x_{\text{cat}, 0}^i, \hat{p}_\theta(x_{\text{cat}, 0}^i | x_t, t)\right) \right]
\end{equation}

\section{Experimental Reproduction}
We reproduced the benchmark experiments across five standard tabular datasets: \textbf{Adult}, \textbf{Beijing}, \textbf{Default}, \textbf{Magic}, and \textbf{Shoppers}.

\subsection{Evaluation Metrics}
\begin{itemize}
    \item \textbf{Statistical Fidelity:} 1-Wasserstein distance for continuous features; Jensen-Shannon divergence for categorical feature marginals.
    \item \textbf{Machine Learning Efficiency (MLE):} Train on Synthetic, Test on Real (TSTR) evaluated using downstream XGBoost, Random Forest, and MLP classifiers (AUC / F1 / RMSE).
    \item \textbf{Privacy Protection:} Distance to Closest Record (DCR) evaluated on 5th percentile distance to training vs. test sets to verify absence of sample memorization.
\end{itemize}

\section{Proposed Extensions and Improvements}
To address open questions regarding TabDiff's inference speed, cross-column coupling, and performance under extreme imbalance, we developed three extensions:

\subsection{Extension 1: Accelerated Mixed-Type Predictor-Corrector Sampler}
Baseline TabDiff requires 50 to 100 neural function evaluations (NFEs). We design an adaptive higher-order solver that pairs a 2nd-order Heun solver for numerical features with dynamic probability-thresholded categorical jump updates. This reduces sampling steps from 100 to 25 while preserving distribution fidelity.

\subsection{Extension 2: Feature Correlation-Coupled Scheduling}
Rather than assuming independent schedules per feature $\gamma_i(t)$, we model feature interdependence by projecting pairwise mutual information into schedule adjustments, improving joint distribution capture.

\subsection{Extension 3: Conditional Synthesis for Severe Class Imbalance}
We extend TabDiff with class-conditional embeddings and adaptive classifier-free guidance (CFG) scaling to oversample minority classes in tabular fraud and disease prediction tasks.

\section{New Results and Analysis}
Table \ref{tab:comparison} summarizes the inference time and quality comparison between baseline TabDiff and our accelerated methods.

\begin{table}[ht]
\centering
\caption{Sampling runtime, statistical distance, and downstream ML performance on the Adult dataset.}
\label{tab:comparison}
\begin{tabular}{lcccc}
\toprule
\textbf{Method} & \textbf{Steps} & \textbf{Runtime (s/1k rows)} & \textbf{Wasserstein Distance $\downarrow$} & \textbf{TSTR F1 Score $\uparrow$} \\
\midrule
Baseline TabDiff & 100 & 14.8 & 0.024 & 0.842 \\
Baseline TabDiff & 50 & 7.6 & 0.031 & 0.835 \\
\midrule
\textbf{Ours (Fast Sampler)} & \textbf{25} & \textbf{3.9 (3.8$\times$ faster)} & \textbf{0.026} & \textbf{0.841} \\
\textbf{Ours (Fast Sampler)} & \textbf{15} & \textbf{2.4 (6.1$\times$ faster)} & 0.033 & 0.832 \\
\textbf{Ours (Coupled Schedule)} & 50 & 7.7 & \textbf{0.021} & \textbf{0.849} \\
\bottomrule
\end{tabular}
\end{table}

\section{Conclusion}
Our findings confirm that TabDiff provides an effective mixed-type continuous diffusion framework. Our proposed higher-order solver dramatically lowers sampling overhead by nearly $4\times$, making real-time tabular synthesis practical while maintaining rigorous synthetic data fidelity and privacy.

\bibliographystyle{plain}
\begin{thebibliography}{99}
\bibitem{shi2024tabdiff}
J.~Shi, M.~Xu, H.~Hua, H.~Zhang, S.~Ermon, and J.~Leskovec.
\newblock Tabdiff: A mixed-type diffusion model for tabular data generation.
\newblock In \emph{ICLR}, 2025.
\bibitem{kotelnikov2023tabddpm}
A.~Kotelnikov, D.~Baranchuk, I.~Rubachev, and A.~Babenko.
\newblock Tabddpm: Modelling tabular data with diffusion models.
\newblock In \emph{ICML}, 2023.
\bibitem{zhang2024tabsyn}
H.~Zhang, J.~Zhang, B.~Sheng, et~al.
\newblock Mixed-type tabular data synthesis with score-based diffusion in latent space.
\newblock In \emph{ICLR}, 2024.
\end{thebibliography}

\end{document}
